\section{Flights, populations and measures}\label{sec:flights}

Which flights, which stretches of flight, and how are they measured?

\paragraph{Declared open flights.} On the FFVL catalogue the pilot declares the type of the flight
(Table~\ref{tab:labels}). A distance, free or with up to three turn points, is open; a triangle, a
quadrilateral or an out-and-return is closed; a hike and fly is neither. We keep the declared open flights:
\FlightsOpen{} of the \FlightsAll{} paraglider flights that depart and are logged every 10~s or finer
(\FlightsClosed{} are closed, \FlightsNeither{} neither).

\begin{table}
\centering\small
\begin{tabular}{@{}l l l r@{}}
\toprule
label & meaning & class & flights \\
\midrule
Dist 3 pts & distance, 3 turn points & open & 40\,094 \\
Dist 2 pts & distance, 2 turn points & open & 25\,399 \\
Dist libre & free distance & open & 694 \\
Dist 1 pt & distance, 1 turn point & open & 658 \\
triangle & triangle & closed & 51\,181 \\
triangle FAI & FAI triangle & closed & 34\,979 \\
Quadrilatère & quadrilateral & closed & 703 \\
Aller-Retour & out and return & closed & 617 \\
Marche et Vol bouclé & hike and fly, closed & closed & 3 \\
Marche et Vol FAI & hike and fly, FAI triangle & closed & 1 \\
Marche et Vol & hike and fly & neither & 269 \\
\bottomrule
\end{tabular}

\caption{The labels of the FFVL catalogue, their class and their flights (paragliders that depart, logged
every 10~s or finer).}
\label{tab:labels}
\end{table}

\paragraph{Groups.} The labels are those of the note \emph{Group flights in free flight}
(\texttt{methods/group-identification}). Two pilots are in contact when one is on the path of the
other within $\delta=2$~min, $\varepsilon_h=500$~m horizontally and $\varepsilon_z=300$~m vertically
(headings compared only in glide) for at least $\tau^*=\tauStar$~min. The configurations of an instant
are the connected components of the contact graph, and a group keeps the same members over a maximal
interval. The groups are those of every flight of the archive, whoever the companions.

\paragraph{Populations.} After its departure, the first fix 3~km from the takeoff, every instant of a
flight is in a group or solo. Three populations:
\begin{itemize}
  \item $\Sl$, \emph{solo flights}: flights that never entered a group, the whole of their time after
        the departure;
  \item $\G$, \emph{group time}: the stretches in a group of the flights that entered one at least once;
  \item $\A$, \emph{the same flights alone}: the solo stretches of those same flights, with the same pilots,
        gliders, loggers and days as $\G$.
\end{itemize}

\paragraph{Measure.} A \emph{piece} is a stretch within one segment of the track, on a 10-s grid of its
clock. For a flight $f$ and a population, the time-averaged mean squared displacement pools the origins of
all its pieces in that population,
\begin{equation}\label{eq:tamsd}
  \overline{\delta^2_f}(\tau)=\frac{1}{N_f(\tau)}\sum_{t}\bigl\|\mathbf r(t+\tau)-\mathbf r(t)\bigr\|^2 ,
\end{equation}
with $t$ and $t+\tau$ in the same piece. The ensemble $\MSD(\tau)$ is its mean over the flights, one weight
per flight. $\Hloc(\tau)$ is half the slope of $\log\MSD$ against $\log\tau$ over the lags within $\pm0.25$
decades of $\tau$: 1 for a straight line, $\tfrac12$ for a random walk. $H$ is half the slope over all the
lags of a curve, from 10~s to its end. Bands are 90\% bands of a two-way bootstrap of pilots and blocks
(\NBoot{} draws of each, Appendix~\ref{app:bootstrap}).

\paragraph{Datasets.} A lag $\tau$ needs pieces of at least $\tau/0.8$~s, so that the origins of every piece
cover the lag in the same share of it. A dataset to the lag $T$ keeps, in each population, its pieces of at
least $T/0.8$~s and measures them at every lag up to $T$: the population does not change along a curve
(Appendix~\ref{app:datasets}). Three datasets:
\begin{itemize}
  \item \emph{1000}: $T=\RawTauOne$~s, pieces of at least \RawShortestOne~s;
  \item \emph{2760}: $T=\RawTauTwo$~s, pieces of at least \RawShortestTwo~s;
  \item \emph{longest}: $T=\RawTauLong$~s, the longest lag at which each population keeps \DeclMinFlights{} flights, pieces of at least \RawShortestLong~s.
\end{itemize}
The quoted values are $\Hk$ on dataset 1000 and $\Hq$ on dataset 2760. The lag of $\Hq$ is 1840~s, and its window of $\pm0.25$ decades
runs from \HqLowest{} to \HqHighest~s. A longer dataset would complete the window, but it keeps fewer pieces, since a piece must last $T/0.8$~s: 2760~s is the
longest lag of the grid at which the comparison on the same route, season, hour and wing (Section~\ref{sec:matched}) still keeps \MatMin{} pieces of $\G$ per level.
On dataset 2760 the window is therefore cut at 2760~s and uses the lags from \HqLowest{} to 2760~s only (Appendix~\ref{app:hloc}); that of $\Hk$ is cut at 1000~s.

\paragraph{Phase statistics.} Flights logged at 1~Hz carry the phases of Vilpellet's method: climbs, searches and glides.
Phases exist on 1-Hz flights only, so the phase statistics rest on fewer flights than the MSD
(\RawPhaseFlightsS{} in $\Sl$, \RawPhaseFlightsA{} in $\A$ and \RawPhaseFlightsG{} in $\G$ against \RawFlightsOneS, \RawFlightsOneA{} and \RawFlightsOneG{} on dataset 1000).
A run belongs to the unit that holds its middle. Seven statistics, each a ratio $n/d$ of two terms summed over the runs of a unit
(Table~\ref{tab:stats}); the phase time $T_u$ of a unit $u$ is the sum of the durations of its runs.

\begin{table}
\centering\footnotesize
\setlength{\tabcolsep}{3pt}
\begin{tabular}{@{}l l l l l@{}}
\toprule
statistic & numerator $n$ & denominator $d$ & runs counted & $r_u$, $q_u$ \\
\midrule
time in climb (\%) & time in climbs & phase time & climbs & $n_u/T_u$, constant \\
time in search (\%) & time in searches & phase time & searches, found and vain & $n_u/T_u$, constant \\
searches ending in a climb (\%) & searches followed by a climb & searches & searches & $n_u/T_u$, $d_u/T_u$ \\
climbs per hour & climbs & phase time (h) & climbs & $n_u/T_u$, constant \\
climb rate (m/s) & height gained in climbs & time in climbs & climbs & $n_u/T_u$, $d_u/T_u$ \\
glide speed (km/h) & straight-line distance of glides & time in glides & glides & $n_u/T_u$, $d_u/T_u$ \\
height lost per km of glide (m/km) & height lost in glides & distance of glides & glides & $n_u/T_u$, $d_u/T_u$ \\
\bottomrule
\end{tabular}
\caption{The seven phase statistics. Unit $u$ of the average: the flight in Section~\ref{sec:raw} (its runs on its pieces of dataset 1000
in the population), the piece or control segment in Sections~\ref{sec:matched} and \ref{sec:decomposition}, the window in Section~\ref{sec:prepost}.}
\label{tab:stats}
\end{table}

\paragraph{Estimator: each unit counts once.} With $w_u$ the weight of a unit (1 for the value, the bootstrap count for a draw),
\begin{equation}\label{eq:estimator}
  r_u=\frac{n_u}{T_u},\qquad q_u=\frac{d_u}{T_u},\qquad
  \text{stat}=\frac{\sum_u w_u\, r_u}{\sum_u w_u\, q_u}.
\end{equation}
For the two time shares and the climbs per hour $q_u$ is constant, and the statistic is the mean of the unit's share or rate. For the
other four it does not divide by the few events of one unit: a unit with one search gives 0\% or 100\%.
In words: per unit, the numerator and the denominator are divided by its phase time, the two are averaged over the units and the statistic is the ratio of the averages.
A unit weighs by its share of phase time in the denominator, not by its length: no long flight dominates, as in the MSD, which has one weight per unit. The ratio of sums, $\sum n_u/\sum d_u$, weights a unit by its phase time
(Appendix~\ref{app:phases}, with an example).

\paragraph{Survival of the phase durations.} For climbs, searches that end in a climb and glides, the survival function
$P(\tau_\text{phase}>\tau)$ of the number of fixes of a run (1~Hz, so seconds), with the duration trimmed as in the
note on phases. Every run weighs 1; its band holds at each $t$ (Appendix~\ref{app:bootstrap})
(Appendix~\ref{app:survival}).
